\chapter{Background}\label{sec:background}

%Hire are Background
%
%Fokus on: 1) give the reader enough technical foundation to understand your method and experiments without consulting other papers, and (2) position your work within the existing literature, making the gap you fill explicit. 
%
%Split into : (a) UED foundations (Underspecified POMDP, Regret as the curriculum objective, Domain Randomization (DR), PAIRED, PLR, ACCEL, SFL,  diffusion models to directly generate high-regret environments
%
%Background for RL: basics of RL auch mitnehmen (ausser ML) + math formel
%
%
%(b) cooperative MARL, 
%- Dec-POMDP: Two paradigms, Independent learning (IPPO) and CTDE methods (MAPPO, QMIX, VDN)
%- varying the partner is not enough; varying the environment is necessary to produce truly robust cooperative agents. This directly motivates UED.
%- Benchmarks that enable (or limit) investigation: Overcooked, OvercookedV2, SMAC/SMACv2, JaxMARL, Craftax-MA.
%
%(c) the intersection: what exists, what doesn't, and why is the gap non-trivial
%
%Related Works + Differenzierung:
%- wie wurde dieses Problen bereits geloest
%- Differenzierung: warum wir machen wenn es bereint gemacht/untersucht wurde

%____________________________________________________________________________________________________

\section{Reinforcement learning}

Reinforcement learning (RL) refers to a subfield of machine learning in which an agent learns through interaction with an environment. Sutton and Barto (2018) \cite{sutton2018reinforcement} define this paradigm as a process in which a sequence of decisions is made to maximise a cumulative reward signal. In contrast to supervised learning, where an external teacher specifies correct target values, in RL the agent receives only a scalar reward signal as feedback on its actions. The agent must therefore independently determine which actions yield the highest return in the long term.  The authors highlight three key characteristics of this approach: Firstly, it is a closed-loop system in which the agent’s actions influence its future inputs. Secondly, the agent receives no explicit instructions as to which actions are optimal. Thirdly, the consequences of actions can extend over extended periods of time.

%
%The mathematical framework of reinforcement learning (RL) is a computational paradigm for automated learning, focused on solving sequential decision-making problems through direct interaction with complex dynamic systems. Unlike classical supervised learning, where the system receives explicit class labels or target values, in the RL paradigm the agent independently optimises its behaviour based on emergent experience and perceptual feedback.

\subsubsection{The basic components of the RL model}
The fundamental architecture of RL is based on the conceptual distinction and closed-loop interaction between two key entities: the Agent and the Environment. According to Sutton and Barto’s \cite{sutton2018reinforcement} fundamental definition, the basic RL setting consists of the following elements:

\begin{itemize}

\item \textbf{Agent}: A control subsystem (controller) capable of perceiving information about the state of the external environment, interpreting it, making decisions regarding the choice of actions, and adapting its own parameters (learning) during operation.

\item \textbf{Environment}: A dynamic system that lies outside the Agent’s immediate internal control but is subject to its influence. The environment encompasses all physical or virtual constraints, transition laws between states, and feedback generation mechanisms.

\item \textbf{State} ($s \in \mathcal{S}$): a complete or partial informational representation of the environment’s configuration at a specific discrete point in time $t$. The state space $\mathcal{S}$ can be either finite (discrete) or infinite-dimensional (continuous). Ideally, a state $s_t \in \mathcal{S}$ satisfies the Markov property, i.e., it contains all relevant retrospective information necessary for predicting the future response of the system.

\item \textbf{Action} ($a \in \mathcal{A}(s)$): A control signal or an act by which the Agent exerts a specific influence on the Environment. The set of available actions $\mathcal{A}(s)$ is a function of the current state, reflecting the dynamic constraints imposed on the Agent’s behaviour in specific system configurations.

\item \textbf{Reward} ($r \in \mathcal{R} \subset \mathbb{R}$): A scalar numerical feedback signal $R_t$ generated by the Environment after each step of the Agent. The physical meaning of the reward lies in the local evaluation of the quality of the transition $(S_{t-1}, A_{t-1}) \to S_t$. The reward mechanism is the sole carrier of information about the global control objective.

\item \textbf{Strategy / Policy} (Policy, $\pi$): A specification of the Agent’s behaviour that defines the mapping from the state space to the action space. In the general stochastic case, it is formalised as a conditional probability distribution:
\[
\pi(a|s) = \mathbb{P}(A_t = a \mid S_t = s), \quad \sum_{a \in \mathcal{A}(s)} \pi(a|s) = 1, \quad \forall s \in \mathcal{S}.
\]
\item \textbf{Episode}: A sequence of interactions that naturally breaks down into time intervals, each ending in a specific state—the terminal state. Problems with this structure are called episodic.

\item \textbf{Return} ($G_t$): A function of the sequence of rewards that the agent seeks to maximise. In the most general form, taking discounting into account, the return is defined as:
\[
G_t \doteq R_{t+1} + \gamma R_{t+2} + \gamma^2 R_{t+3} + \dots
= \sum_{k=0}^{\infty} \gamma^k R_{t+k+1}
\]

\item \textbf{Discount factor} ($\gamma$):A parameter in the interval $0 \le \gamma \le 1$ that determines the present value of future rewards. When $\gamma < 1$, the sum in the revenue formula is finite if the reward sequence is bounded.

\end{itemize}

\subsubsection{The agent-environment interaction cycle}

The agent’s interaction with the environment occurs continuously through a sequence of discrete time steps $t = 0, 1, 2, 3, \dots$. This iterative process follows the pattern shown in Figure~\ref{fig:imgs/RL_Loop}: Agent $\rightarrow$ Action $\rightarrow$ Environment $\rightarrow$ Reward + Next state. At each time step $t$, the agent perceives the state of the environment $S_t$ and, based on this, performs action $A_t$. At the next time step, as a partial consequence of this action, the agent receives a numerical reward $R_{t+1}$ from the environment and finds itself in a new state $S_{t+1}$. This cycle forms a continuous trajectory of the agent’s experience.
\begin{figure}[htbp]
    \centering
    \includegraphics[width=0.8\textwidth]{imgs/RL_Loop.png}
    \caption{Agent-environment interaction cycle}
    \label{fig:imgs/RL_Loop}
\end{figure}

\subsubsection{Formal Objective}

The entire paradigm of reinforcement learning is based on the so-called ``\emph{reward hypothesis}'': any goal or objective can be formalised as the maximisation of the expected value of the cumulative sum of the received scalar reward signal. To formalise this long-term goal, the concept of return ($G_t$) is introduced. For continuous-time problems, the return is defined as the sum of discounted future rewards:

\begin{equation}
G_t = \sum_{k=0}^{\infty} \gamma^k R_{t+k+1}
\label{eq:return}
\end{equation}

To assess how beneficial it is for an agent to be in a particular state $s$ when following a given policy $\pi$, the state-value function ($v_{\pi}(s)$) is used. It is defined as the expected return starting from state $s$:
\begin{equation}
v_{\pi}(s)
=
\mathbb{E}_{\pi}\!\left[G_t \mid S_t = s\right]
=
\mathbb{E}_{\pi}\!\left[
\sum_{k=0}^{\infty}
\gamma^k R_{t+k+1}
\;\middle|\;
S_t = s
\right]
\end{equation}


Thus, the formal objective of an agent in reinforcement learning is to find the optimal policy ($\pi^{*}$), whose expected reward is greater than or equal to that of any other policy for every state without exception. Mathematically, this means achieving the optimal state value function ($v^{*}(s)$), defined as:
\begin{equation}
v^{*}(s) = \max_{\pi} v_{\pi}(s)
\end{equation}
which maximises the expected reward for every state of the environment.

Thus, rather than globally optimising an abstract function, RL algorithms seek behaviour that will maximise the expected return ($v^{*}(s)$) for each local state in which the agent may find itself.


\subsubsection{Bellman Equations}

A fundamental property of value functions is their recursive structure, which is expressed by the Bellman equation. For the state-value function, the following relationship holds as formulated by Sutton and Barto (2018):
\[
v_\pi(s) = \sum_{a} \pi(a \mid s) \sum_{s', r} p(s', r \mid s, a) \left[r + \gamma v_\pi(s')\right].
\]

This equation expresses that the value of a state equals the expected discounted value of successor states plus the expected immediate reward. It forms the foundation of all dynamic programming and reinforcement learning algorithms.

For the optimal state-value function \(v^*(s)\), the Bellman optimality equation is given by:
\[
v^*(s) = \max_{a} \sum_{s', r} p(s', r \mid s, a) \left[r + \gamma v^*(s')\right].
\]

Similarly, for the optimal action-value function \(q^*(s, a)\), it holds that:
\[
q^*(s, a) = \sum_{s', r} p(s', r \mid s, a) \left[r + \gamma \max_{a'} q^*(s', a')\right].
\]

Sutton and Barto \cite{sutton2018reinforcement} show that any policy that selects, in every state, only the action with the highest \(q^*\)-value (i.e., greedy with respect to \(v^*\) or \(q^*\)) constitutes an optimal policy \(\pi^*\).


\subsubsection{Markov Decision Processes}

The mathematical foundation of Reinforcement Learning is the Markov Decision Process (\textit{MDP}). Following the standard work of Sutton and Barto \cite{sutton2018reinforcement}, a finite MDP is defined by the tuple \((S, A, p, r, \gamma)\), where \(S\) denotes the finite set of states, \(A\) the finite set of actions, \(p\) the transition dynamics, \(r\) the reward function, and \(\gamma \in [0,1]\) the discount factor.

At each discrete time step \(t\), the agent is in a state \(S_t \in S\) and selects an action \(A_t \in A(S_t)\). The environment then transitions to a new state \(S_{t+1}\), and the agent receives a numerical reward \(R_{t+1}\).

The dynamics of the MDP are fully described by the joint transition and reward probability:
\[
p(s', r \mid s, a)
=
\Pr\{S_{t+1} = s', R_{t+1} = r \mid S_t = s, A_t = a\}.
\]

This formulation assumes the Markov property. As explained by Sutton and Barto (2018) \cite{sutton2018reinforcement}, the probability of a state transition depends solely on the current state and the current action, rather than on the entire history. Formally, this means that the complete dynamics are identical to the conditional probability given the full history
\[
S_0, A_0, R_1, \dots, S_t, A_t.
\]


\subsubsection{Value Functions}

Almost all reinforcement learning algorithms involve estimating value functions—functions of states (or state–action pairs) that assess how good it is for an agent to be in a given state or how good it is to take a certain action in that state. The notion of ``how good'' is defined in terms of future rewards that the agent is expected to receive, or more precisely, in terms of the expected return. Since expected rewards depend on the actions chosen by the agent, value functions are always defined relative to a particular policy \(\pi\).

In reinforcement learning theory, two main types of value functions are distinguished:

1. **State-value function** \(v_\pi(s)\)  
The value of a state \(s\) under a policy \(\pi\), denoted \(v_\pi(s)\), is the expected return that an agent obtains when starting from state \(s\) and thereafter following policy \(\pi\). Formally, for Markov Decision Processes (MDPs), this function is defined as follows:

\[
v_\pi(s) = \mathbb{E}_\pi \left[ G_t \mid S_t = s \right]
= \mathbb{E}_\pi \left[ \sum_{k=0}^{\infty} \gamma^k R_{t+k+1} \mid S_t = s \right]
\]

This mathematical construct is called the state-value function for policy \(\pi\). By definition, the value of any terminal state is always equal to zero.

2. **Action-value function** \(q_\pi(s,a)\)  
Similarly, the value of taking a specific action \(a\) in state \(s\) under policy \(\pi\), denoted \(q_\pi(s,a)\), is defined. It is the expected return when the agent starts from state \(s\), takes action \(a\), and then continues to follow policy \(\pi\). Mathematically, it is written as:

\[
q_\pi(s,a) = \mathbb{E}_\pi \left[ G_t \mid S_t = s, A_t = a \right]
= \mathbb{E}_\pi \left[ \sum_{k=0}^{\infty} \gamma^k R_{t+k+1} \mid S_t = s, A_t = a \right]
\]

This function is called the action-value function for policy \(\pi\).

\subsubsection{Dynamic Programming}

Dynamic Programming (DP) comprises a class of algorithms that can compute optimal policies given a known environment model. The framework of Generalized Policy Iteration (GPI), as described by Sutton and Barto (2018), alternates between two central procedures: policy evaluation and policy improvement.

During policy evaluation, the Bellman equation is iteratively solved to approximate \(v_\pi\) for a given policy \(\pi\). During policy improvement, a better policy is constructed by selecting, in each state, the action that maximizes the value function. The repeated application of these two steps converges to the optimal policy.

However, DP methods assume full knowledge of the environment dynamics \(p(s', r \mid s, a)\) and become computationally expensive in large state spaces.

\subsubsection{Monte Carlo Methods and Temporal-Difference Learning}

In contrast to dynamic programming methods, Monte Carlo (MC) methods do not require a model of the environment. They estimate value functions based on complete episodes by using the actual return \(G_t\) as a target:
\[
V(S_t) \leftarrow V(S_t) + \alpha \left[ G_t - V(S_t) \right].
\]

Temporal-Difference (TD) learning, extensively discussed by Sutton and Barto (2018), combines properties of both approaches. Like Monte Carlo methods, TD learning learns directly from experience without requiring a model; like dynamic programming, it updates estimates based on other learned estimates (bootstrapping).

The basic TD(0) update rule is given by:
\[
V(S_t) \leftarrow V(S_t) + \alpha \left[ R_{t+1} + \gamma V(S_{t+1}) - V(S_t) \right].
\]

The term
\[
\delta_t = R_{t+1} + \gamma V(S_{t+1}) - V(S_t)
\]
is known as the temporal-difference error (TD error) and represents a central concept in reinforcement learning.

\subsubsection{TD Control Methods: SARSA and Q-Learning}

For the control problem, i.e., finding an optimal policy, temporal-difference methods are applied to action-value functions. Two central algorithms in this context are SARSA (on-policy) and Q-learning (off-policy).

SARSA updates the action-value function along the trajectory of the policy that is actually being followed:
\[
Q(S_t, A_t) \leftarrow Q(S_t, A_t) + \alpha \left[ R_{t+1} + \gamma Q(S_{t+1}, A_{t+1}) - Q(S_t, A_t) \right].
\]

The name SARSA is derived, as described by the authors, from the quintuple \((S_t, A_t, R_{t+1}, S_{t+1}, A_{t+1})\) used in each update.

In contrast, Q-learning directly approximates the optimal action-value function \(q^*\), independent of the policy being followed:
\[
Q(S_t, A_t) \leftarrow Q(S_t, A_t) + \alpha \left[ R_{t+1} + \gamma \max_{a} Q(S_{t+1}, a) - Q(S_t, A_t) \right].
\]

The key difference is that Q-learning uses the maximum action value over all possible actions in the next state (\(\max_a\)), whereas SARSA uses the value of the actually executed next action. Under mild conditions, Q-learning converges to \(q^*\) with probability 1.

\subsubsection{Exploration and Exploitation}

A central dilemma in reinforcement learning is the trade-off between exploration (trying new actions) and exploitation (selecting actions known to yield high rewards). Sutton and Barto (2018) discuss several approaches to address this problem.

A widely used strategy is the \(\varepsilon\)-greedy policy, where the agent selects the currently best-known action with probability \(1 - \varepsilon\) and a random action with probability \(\varepsilon\).

Other approaches described by the authors include the Upper Confidence Bound (UCB) method and softmax-based action selection strategies.


---


The research of robust and cooperative agens in multi-agent systems requires a systhesis of different  research fields: Unsupervised Environment Design (UED), multi-agent reinforcement learning (MARL) and automated generation of learning plans.


\subsection*{Unsupervised Environment Design}
Unsupervised environment design was introduced as an approach to automatically generate learning environments that enable agents to develop more robust behaviour. One of the seminal works is by \cite{baker2020}, who, in \emph{‘Emergent Tool Use from Multi-Agent Autocurricula’}, demonstrate that complex skills can emerge through adversarial environment generation. In this process, a ‘teacher’ generates environments, while a ‘student’ attempts to solve them.

Building on this, the UED paradigm has been systematised, for example through work such as PAIRED (Protagonist Antagonist Induced Regret Environment Design), which specifically generates challenging but solvable environments (\cite{dennis2020}). This work demonstrates that automatically generated curricula enable greater generalisation than static training environments. However, much of the research to date has focused on single-agent scenarios, meaning that its applicability to cooperative multi-agent systems has only been investigated to a limited extent so far.

UED is often described as a two-player game between an level-solving agent and a level-generating adversary \cite{monette2025optimisationframeworkunsupervisedenvironment}. A key metric for level selection is regret \cite{dennis2020}– the difference between the potentially optimal performance on a level and the actual performance of the current agent. The best-known methods can be divided into two categories:
\begin{itemize}
    \item Replay-based UED methods: Rather than having a neural network generate environments from scratch each time, algorithms such as Prioritised Level Replay (PLR) \cite{JiangPLR2021} rely on a buffer of previously seen levels. PLR evaluates randomly generated levels based on their learning potential (e.g. approximated by the Positive Value Loss) and trains the agent with a priority on those environments that exhibit the highest approximated regret, leading to significantly more robust policies in grid-based navigation tasks \cite{parkerholder2023ACCEL}.

    \item Evolutionary and mutation-based UED methods: The ACCEL algorithm (Adversarially Compounding Complexity by Editing Levels) \cite{parkerholder2023ACCEL} extends PLR with a mutation operator. It makes small adjustments to levels that already exhibit high regret, thereby continuously pushing the agent to the limits of its capabilities, which causes the complexity of the environments to increase steadily

    \item Learnability-based approaches: Recent research shows that regret approximations are often inaccurate in practice and do not always correlate with the true learning potential. Sampling For Learnability (SFL) addresses this by measuring the learnability of a level based on the variance of the success rate ($p(1-p)$) \cite{rutherford2024sfl}. SFL consequently prioritises environments that the agent can solve ‘sometimes, but not always’, leading to significantly more robust policies in grid-based navigation tasks.

\end{itemize}




%_________________________________

\subsubsection{UED in multi-agent scenarios}

While UED has been extensively researched in a single-agent context, its application to multi-agent systems is still in its early stages\cite{dennis2020, JiangPLR2021, parkerholder2023ACCEL}, . In multi-agent domains, there is the additional challenge that the complexity of a task depends not only on the static environment, but is also significantly determined by the strategies of the other players (co-players) \cite{samvelyan2023maestr}

Previous approaches to generating auto-curricula in MARL have primarily focused on the evolution of opponent policies in competitive games, such as through fictitious self-play (FSP) or population-based training \cite{samvelyan2023maestr}. The recently introduced MAESTRO (Multi-Agent Environment Design Strategist for Open-Ended Learning) algorithm is one of the first approaches to transfer UED to competitive MARL environments (two-player zero-sum games) \cite{samvelyan2023maestr}. MAESTRO combines PLR with population-based learning to generate a common curriculum across environments and opponent policies, thereby making agents more robust to both environmental changes and unknown opponents \cite{samvelyan2023maestr}

However, in the cooperative setting that is the focus of this work, there are hardly any comparable approaches \cite{ruhdorfer2025overcookedgeneralisationchallengeevaluating}. Much of the existing research on multi-agent UED either does not address fully cooperative settings, focuses on rudimentary design spaces, or investigates pathfinding without direct agent interaction \cite{ruhdorfer2025overcookedgeneralisationchallengeevaluating}. Recent initiatives such as the Overcooked Generalisation Challenge \cite{ruhdorfer2025overcookedgeneralisationchallengeevaluating} or the Craftax-MA \cite{} benchmark clearly demonstrate that classical MARL algorithms perform significantly worse in procedurally generated, open-ended cooperative tasks with long time horizons and heterogeneous requirements \cite{omari2025MACraftax}. Early frameworks in the ‘Overcooked’ domain have already defined Dec-UPOMDPs (Underspecified POMDPs with Shared Rewards), in which layout, object distributions and starting positions are parameterised to enable zero-shot coordination across different environments \cite{ruhdorfer2025overcookedgeneralisationchallengeevaluating}

%____________________________

To overcome the poor generalisation ability of RL agents, they must be trained across a wide variety of environments. A classic approach to this is domain randomisation (DR) \cite{tobin2017domainrandomizationtransferringdeep}, in which environment parameters are randomly sampled. However, as DR often generates unsolvable or trivial environments \cite{dennis2020}, Unsupervised Environment Design (UED) was introduced as a scalable approach, in which a curriculum is adaptively adjusted to the agent’s current capabilities \cite{ruhdorfer2025overcookedgeneralisationchallengeevaluating}





\subsection*{Multi-Agent Reinforcement Learning and Cooperation}
In the field of multi-agent reinforcement learning, the emergence of cooperative behaviour presents a key challenge. Traditional approaches such as Independent Q-Learning~\cite{TanMARL1993} suffer from non-stationarity, as the policies of other agents change during training. Advanced methods such as QMIX~\cite{RashidQMIX2018}, VDN (Value Decomposition Networks)~\cite{SunehagVDN2018}, and MADDPG~\cite{Lowe2017} address this problem through centralised training mechanisms with decentralised execution. For cooperative scenarios, the Overcooked\footnote{\url{https://github.com/HumanCompatibleAI/overcooked_ai}} domain~\cite{Carroll2019} was established in particular, in which agents must coordinate tasks together. This environment allows for a controlled investigation of cooperation, communication and coordination strategies.

%____________________________

Solving complex coordination problems involving multiple agents is a key area of research in artificial intelligence. The fully cooperative MARL problem is often formally modelled as a decentralised partially observable Markov decision process (Dec-POMDP) \cite{omari2025MACraftax}, in which all agents act on the basis of their local observations but optimise a shared, scalar team reward signal \cite{dewitt2020independentlearningneedstarcraft}


A straightforward approach is Independent Learning (IL), in which each agent treats its environment as stationary and regards the other agents merely as part of that environment \cite{dewitt2020independentlearningneedstarcraft}. However, as this leads to learning instabilities, the Centralised Training with Decentralised Execution (CTDE) \cite{Lowe2017} paradigm has become established
. Here, training is carried out centrally in a simulation where global state information is available for learning a common value function (critic), while the execution of the learned policies is strictly decentralised and based on local observations \cite{Sartoretti_2019}

Although MARL algorithms such as QMIX \cite{RashidQMIX2018} or COMA \cite{foerster2024counterfactualmultiagentpolicygradients}achieve high performance levels in standard benchmarks (e.g. SMAC \cite{samvelyan2019starcraftmultiagentchallenge}), they typically train agents using so-called self-play (SP) \cite{Gupta2017}. As a result, agents are heavily reliant on the conventions established during training. When these agents have to interact in new environments or with unknown partners (a problem known as zero-shot coordination or ZSC \cite{hu2021otherplayzeroshotcoordination}), they often fail to generalise effectively
.


\subsection*{Autocurricula and Level-Generation}

Another relevant area of research is automatic curricula that adaptively select training data or environments.
Key approaches include: Prioritized Level Replay (PLR) (\cite{JiangPLR2021}) which selects training levels based on learning progress and difficulty; ACCEL (Adaptive Curriculum via Contrastive Learning): Uses contrastive learning to generate meaningful training progressions. SFL (Sampling For Learnability) is a simple approach that directly chooses levels that optimise learnability alongside randomly generated levels. \cite{rutherford2024sfl}. MAESTRO (Open-Ended Learning Framework) \cite{samvelyan2023maestr}Combines evolution and curriculum learning for the open-ended generation of tasks.

These methods demonstrate that dynamic curricula are crucial for efficient learning. Nevertheless, they have predominantly been evaluated in single-agent contexts.

\subsection*{Combination of UED and MARL}

The integration of UED into cooperative multi-agent scenarios has so far been insufficiently researched. Although UED has been successfully deployed in competitive environments, cooperative settings give rise to additional requirements, particularly with regard to coordination and the joint achievement of goals.

Parameterised environments — such as varying layouts, resource distributions or partial observability — offer the possibility of specifically modulating cooperation requirements. In this context, the question arises as to whether principles such as learnability, as defined in SFL, are also transferable to multi-agent scenarios, for example through shared success metrics or coordination-dependent levels of difficulty




%\subsection{Policy, Return, and Value Functions}

%A policy \(\pi\) describes a mapping from states to probability distributions over actions. For a stochastic policy, \(\pi(a \mid s) = \Pr\{A_t = a \mid S_t = s\}\). The primary objective of the agent, following Sutton and Barto (2018), is to maximize the expected return \(G_t\), which is defined as the discounted sum of future rewards:
%\[
%G_t = R_{t+1} + \gamma R_{t+2} + \gamma^2 R_{t+3} + \dots = \sum_{k=0}^{\infty} \gamma^k R_{t+k+1}.
%\]

%The discount factor \(\gamma\) determines the weight of future rewards. For \(\gamma = 0\), the agent maximizes only immediate reward, whereas for \(\gamma \to 1\), future rewards become increasingly important.

%To formally evaluate the performance of a policy, the authors introduce two fundamental value functions. The state-value function \(v_\pi(s)\) gives the expected return when the agent starts in state \(s\) and subsequently follows policy \(\pi\):
%\[
%v_\pi(s) = \mathbb{E}_\pi \left[ G_t \mid S_t = s \right].
%\]

%Similarly, the action-value function \(q_\pi(s, a)\) defines the expected return when action \(a\) is taken in state \(s\) and the agent thereafter follows policy \(\pi\):
%\[
%q_\pi(s, a) = \mathbb{E}_\pi \left[ G_t \mid S_t = s, A_t = a \right].
%\]



\subsection{Multi-agent reinforcement learning}

\subsubsection{Dec-POMDPs.}
Dec-POMDPs. A Dec-POMDP is defined as a tuple $(n, S, A, T, O, O, R, \gamma)$ \cite{oliehoek2008optimal}. Here, $n \in \mathbb{N}$ denotes the number of agents, $S$ the state space, $A$ is the action space for each agent, and $O$ represents an agent’s observation space. The transition probability function $T : S \times A^n \rightarrow \Delta(S)$ dictates the evolution of the environment, the observation function $O : S \times A^n \rightarrow \Delta(O^n)$ determines what agents can perceive, and $R : S \times A^n \rightarrow \mathbb{R}$ is the reward function. Finally, $\gamma \in [0,1)$ is the discount factor. At every timestep $t$, each agent $i$ selects an action $a_t^i \in A$. The environment then transitions from state $s_t \in S$ to $s_{t+1} \in S$ according to $T(s_t, a)$ with probability $P(s_{t+1} \mid s_t, a_t)$, where $a_t = (a_t^1, \ldots, a_t^n)$ represents the joint action. A reward $r_t \in \mathbb{R}$ is drawn based on $R(s_t, a_t)$, and each agent receives an observation $o_{t+1}^i$, with the joint observation $o_{t+1} = (o_{t+1}^1, \ldots, o_{t+1}^n)$ sampled from $O(s_{t+1}, a_t)$. The primary objective is to learn policies $\pi^i : (O^i \times A^i)^* \rightarrow \Delta(A^i)$ for each agent $i$, which maximises the expected cumulative discounted reward $\mathbb{E}\left[\sum_{t=0}^{T} \gamma^t r_t \right]$.

\subsubsection{Partial Observability and Stochasticity}


\subsection{Methods for UED}
\begin{itemize}
    \item Constant: domain randomization \cite{}
    \item Regret-Based:  \cite{dennis2020} minimax regret, \cite{parkerholder2023ACCEL} PLR, \cite{JiangPLR2021}
    
\end{itemize}

\subsubsection{Sparse and dense reward}


\section{Related works}